% !Mode:: "TeX:UTF-8"

\chapter{控制算法设计}

\section{引言}

本章在前述运动学建模
与位置控制基础验证的
基础上，进一步面向样
机系统的实际实现需求，
建立单电机闭环力矩控
制与四电机协同控制算
法。对于本文所研究的
高动态连续旋转平面四
绳绳驱并联机器人而言，
单纯依赖位置指令驱动
难以直接保证末端平台
在运动过程中的受力品
质与加速度连续性。因
此，有必要从单执行单
元的力矩控制出发，逐
步构建末端位姿误差补
偿、广义力修正以及张
力映射的完整控制链路。

本章的核心目标包括三
个方面。其一，针对单
电机执行单元建立闭环
力矩控制方法，在 PID
控制基础上引入摩擦补
偿与扰动观测器，以提
高力矩跟踪精度和抗扰
能力。其二，基于正运
动学反馈构建末端位姿
闭环补偿方法，将末端
轨迹误差转换为平台广
义力修正量。其三，结
合平面四绳冗余驱动机
构的特点，建立广义力
修正到各绳张力增量及
电机目标力矩之间的映
射关系，为后续张力分
配分析和实验验证提供
控制算法基础。

\section{控制系统总体架构}

结合本文样机结构与控
制任务特点，系统控制
架构采用“单电机内环
力矩控制+末端外环位姿
补偿”的双层形式。其
中，内环面向单个执行
单元，重点解决电机输
出力矩的跟踪精度、摩
擦扰动抑制以及动态响
应速度问题；外环面向
末端平台，利用正运动
学重构得到的末端位姿
信息构造误差补偿量，
再通过广义力与张力映
射关系分配至四个电机。

对于第 $i$ 个电机，设
其角位移、角速度和输
入力矩分别为

\begin{equation}
  \phi_i,\qquad
  \omega_i=\dot{\phi}_i,\qquad
  \tau_{m,i}.
\end{equation}

设电机等效转动惯量为
$J_i$，等效粘性阻尼为
$B_i$，卷筒半径为 $r_m$，
绳索张力为 $f_i$，则单
电机执行单元的等效动
力学可写为

\begin{equation}
  J_i\dot{\omega}_i
  =
  \tau_{m,i}
  -
  r_m f_i
  -
  B_i\omega_i
  -
  \tau_{f,i}
  -
  \tau_{d,i},
  \label{eq:ch4-motor-dyn}
\end{equation}

其中，$\tau_{f,i}$ 表示摩
擦力矩，$\tau_{d,i}$ 表示
建模误差、负载扰动及
其他未建模因素引起的
等效扰动力矩。

在系统总体层面，末端
期望轨迹

\begin{equation}
  \bm q_d=
  \begin{bmatrix}
    x_d & y_d & \theta_d
  \end{bmatrix}^{\mathrm T}
\end{equation}

经逆运动学和状态反馈
后，可由正运动学得到
当前末端位姿估计值

\begin{equation}
  \hat{\bm q}=
  \begin{bmatrix}
    \hat{x} & \hat{y} & \hat{\theta}
  \end{bmatrix}^{\mathrm T}.
\end{equation}

进而构造末端位姿误差

\begin{equation}
  \bm e_q=\bm q_d-\hat{\bm q}.
  \label{eq:ch4-eq}
\end{equation}

该误差经外环补偿器处
理后形成末端广义力修
正量，再由张力映射模
块转换为四根绳的张力
增量和相应电机目标力
矩，最终由四个单电机
内环分别完成跟踪。因
此，系统控制信息流可
概括为“轨迹规划—位
姿反馈—误差补偿—张
力映射—电机力矩闭环”。

\begin{figure}[htbp]
  \centering
  \fbox{\begin{minipage}[c][5.0cm][c]{0.84\textwidth}
    \centering
    占位图：控制系统总体架构图。\\
    建议包含末端期望轨迹、逆运动学、\\
    正运动学反馈、末端位姿补偿模块、\\
    张力映射模块以及四个单电机内环。
  \end{minipage}}
  \caption{控制系统总体架构图}
  \label{fig:ch4-control-architecture}
\end{figure}

如\cref{fig:ch4-control-architecture}
所示，本章控制算法并
不直接替代后续张力分
配优化，而是为其提供
可执行的控制框架。其
中，内环负责提高单电
机力矩控制品质，外环
负责将末端位姿误差转
换为广义力修正需求，
而最终张力分配与冗余
调节将在下一章中结合
动力学模型进一步展开。

\section{单电机闭环力矩控制方法}

\subsection{控制对象与问题描述}

对于绳驱并联机器人而
言，各电机通过卷筒收
放绳索间接调节绳张力，
进而影响末端平台所受
广义力。由于四电机协
同控制建立在各执行单
元具备基本稳定性和良
好跟踪性能的前提下，
因此有必要先对单电机
闭环力矩控制方法进行
设计。

定义第 $i$ 个执行单元
的目标张力为 $f_i^\ast$，
其对应的目标卷筒力矩
为

\begin{equation}
  \tau_i^\ast=r_m f_i^\ast.
  \label{eq:ch4-tau-ref-basic}
\end{equation}

设第 $i$ 个执行单元的
实际输出等效力矩为
$\tau_i$，则力矩跟踪误差
定义为

\begin{equation}
  e_{\tau,i}=\tau_i^\ast-\tau_i.
  \label{eq:ch4-etau}
\end{equation}

单电机力矩控制的主要
目标是使 $e_{\tau,i}$ 尽可
能趋近于零，并在存在
摩擦、负载变化及模型
不确定性的条件下保持
较好的跟踪精度与动态
响应。考虑到样机系统
的工程实现难度与实时
控制需求，本文采用
PID 控制、摩擦补偿与
扰动观测器相结合的方
法构建单电机控制器。

\begin{figure}[htbp]
  \centering
  \fbox{\begin{minipage}[c][4.8cm][c]{0.84\textwidth}
    \centering
    占位图：单电机闭环力矩控制对象示意图。\\
    建议给出电机、卷筒、绳索、张力方向、\\
    传感器位置及负载作用关系。
  \end{minipage}}
  \caption{单电机闭环力矩控制对象示意图}
  \label{fig:ch4-single-motor-object}
\end{figure}

\subsection{基于PID的力矩闭环控制}

PID 控制器具有结构简
单、参数物理意义明确
且工程应用成熟等优点，
适合作为单电机力矩闭
环控制的基础。对第
$i$ 个电机，PID 控制输
出写为

\begin{equation}
  \tau_{pid,i}
  =
  K_{p,i}e_{\tau,i}
  +
  K_{i,i}\int e_{\tau,i}\mathrm dt
  +
  K_{d,i}\dot e_{\tau,i},
  \label{eq:ch4-pid}
\end{equation}

其中，$K_{p,i}$、$K_{i,i}$
和 $K_{d,i}$ 分别为比例、
积分和微分增益。

比例项用于根据当前力
矩误差迅速生成校正作
用，积分项用于减小稳
态误差，微分项则用于
改善系统瞬态响应并抑
制误差变化过快导致的
振荡。考虑到实际力矩
信号可能存在高频噪声，
微分项在实现时通常采
用一阶滤波形式，以避
免噪声放大。

仅采用 PID 控制时，
若系统模型较准确、摩
擦影响较小且外部扰动
较弱，则能够在一定程
度上完成力矩跟踪。但
对于本文样机而言，绳
索绕线摩擦、卷筒回差、
负载变化以及机构耦合
都会对力矩跟踪质量产
生影响，因此还需进一
步引入补偿环节。

\begin{figure}[htbp]
  \centering
  \fbox{\begin{minipage}[c][4.8cm][c]{0.84\textwidth}
    \centering
    占位图：基于PID的单电机力矩闭环框图。\\
    建议包含目标力矩、误差计算、PID 控制器、\\
    电机执行对象及实际力矩反馈。
  \end{minipage}}
  \caption{基于PID的单电机力矩闭环框图}
  \label{fig:ch4-pid-block}
\end{figure}

\subsection{摩擦补偿设计}

在绳驱系统中，摩擦主
要来源于电机轴承、减
速机构、卷筒与绳索接
触以及导向部件等环节。
这些摩擦会导致低速段
响应迟滞、换向时死区
加重以及力矩输出不对
称等问题，从而降低控
制精度。为提高单电机
力矩跟踪性能，本文在
PID 控制基础上引入摩
擦补偿项。

考虑工程实现的简洁性，
本文采用库仑摩擦与粘
性摩擦组合的近似模型，
将第 $i$ 个执行单元的
摩擦补偿力矩写为

\begin{equation}
  \hat{\tau}_{f,i}
  =
  \tau_{c,i}\,
  \operatorname{sat}\!\left(
  \frac{\omega_i}{\omega_0}
  \right)
  +
  b_i\omega_i,
  \label{eq:ch4-friction}
\end{equation}

其中，$\tau_{c,i}$ 为等效
库仑摩擦力矩，$b_i$ 为
等效粘性摩擦系数，
$\omega_0$ 为小正数，用
于在低速区域平滑逼近
符号函数，
$\operatorname{sat}(\cdot)$
为饱和函数。

摩擦补偿的物理含义在
于：当电机转速较高时，
补偿项近似为库仑摩擦
与粘性摩擦之和；当电
机接近零速时，补偿项
通过平滑函数避免突变，
从而降低换向抖振。由
于本文更关注工程可实
现性，因此并未进一步
引入更复杂的 LuGre 等
动态摩擦模型，而是采
用参数较少、便于整定
的静态补偿形式。

\begin{figure}[htbp]
  \centering
  \fbox{\begin{minipage}[c][4.6cm][c]{0.84\textwidth}
    \centering
    占位图：摩擦补偿特性曲线示意图。\\
    建议给出速度与摩擦补偿力矩之间的关系，\\
    并标明低速平滑区和高速区特性。
  \end{minipage}}
  \caption{摩擦补偿特性示意图}
  \label{fig:ch4-friction-curve}
\end{figure}

\subsection{扰动观测器设计}

除摩擦外，系统还受到
负载波动、参数失配及
未建模动态等因素影响。
若仅靠固定参数的反馈
控制器，难以在不同工
况下始终获得稳定一致
的控制效果。为此，本
文在单电机力矩闭环中
引入扰动观测器，对总
扰动力矩进行在线估计
与补偿。

将\cref{eq:ch4-motor-dyn}
改写为名义模型形式，

\begin{equation}
  J_{n,i}\dot{\omega}_i
  =
  \tau_{m,i}
  -
  \tau_{dis,i},
  \label{eq:ch4-nominal}
\end{equation}

其中，$J_{n,i}$ 为第 $i$
个电机的名义转动惯量，
$\tau_{dis,i}$ 为总扰动力
矩，综合包含真实惯量
与名义惯量差异、负载
变化、摩擦误差及其他
未建模因素。

基于名义模型，可构造
扰动观测器估计量

\begin{equation}
  \hat{\tau}_{dis,i}
  =
  Q_i(s)
  \left(
  \tau_{m,i}
  -
  J_{n,i}s\omega_i
  \right),
  \label{eq:ch4-dob}
\end{equation}

其中，$Q_i(s)$ 为低通滤
波器，其作用是在抑制
测量噪声的同时估计低
频主导扰动。典型情况
下可取

\begin{equation}
  Q_i(s)=\frac{g_i}{s+g_i},
  \label{eq:ch4-qfilter}
\end{equation}

式中，$g_i$ 为滤波器带
宽。带宽越大，扰动估
计响应越快，但对噪声
也越敏感；带宽越小，
则估计更平滑，但动态
补偿能力下降。因此，
$g_i$ 需要结合执行器响
应速度和采样噪声水平
进行折中选取。

综合 PID、摩擦补偿与
扰动观测器后，第 $i$
个电机的最终控制输入
可写为

\begin{equation}
  \tau_{m,i}^\ast
  =
  \tau_{pid,i}
  +
  \hat{\tau}_{f,i}
  +
  \hat{\tau}_{dis,i}.
  \label{eq:ch4-final-inner}
\end{equation}

上述控制律的物理意义
是：PID 环节负责构造
基本闭环性能，摩擦补
偿环节负责消除主要静
态非线性影响，扰动观
测器则进一步对剩余不
确定项进行实时修正。
这种组合方式兼顾了工
程实现简洁性和控制效
果，是本文单电机闭环
力矩控制的基本形式。

\begin{figure}[htbp]
  \centering
  \fbox{\begin{minipage}[c][5.0cm][c]{0.84\textwidth}
    \centering
    占位图：含摩擦补偿与扰动观测器的单电机\\
    闭环力矩控制框图。建议包含 PID、摩擦补偿、\\
    扰动观测器、被控对象及反馈信号通路。
  \end{minipage}}
  \caption{单电机闭环力矩控制完整框图}
  \label{fig:ch4-inner-loop-full}
\end{figure}

\section{基于正运动学反馈的末端位姿闭环补偿方法}

\subsection{末端位姿误差定义}

在四电机协同控制中，
单电机内环只负责保证
各执行单元具备较好的
力矩跟踪能力，而末端
平台能否按照期望轨迹
运行，还需要依赖任务
空间层面的误差补偿。
因此，本文基于正运动
学反馈构建末端位姿闭
环补偿方法。

根据第 3 章建立的正运
动学求解器，可由实时
绳长测量值反求末端位
姿估计值 $\hat{\bm q}$。在
此基础上，定义末端位
姿误差为

\begin{equation}
  \bm e_q
  =
  \bm q_d-\hat{\bm q}
  =
  \begin{bmatrix}
    e_x & e_y & e_\theta
  \end{bmatrix}^{\mathrm T}.
  \label{eq:ch4-task-error}
\end{equation}

相应地，末端位姿误差
变化率定义为

\begin{equation}
  \dot{\bm e}_q
  =
  \dot{\bm q}_d-\dot{\hat{\bm q}}.
  \label{eq:ch4-task-edot}
\end{equation}

上述误差量分别反映了
末端平台在平移与转动
方向上的偏差，是外环
补偿器的直接输入。与
单纯通过电机编码器比
较位置误差不同，该误
差定义建立在正运动学
重构基础上，更能反映
末端平台真实运动状态，
也便于将后续补偿量直
接构造为末端平台广义
力修正需求。

\subsection{任务空间误差到广义力修正量的构造}

对于平面四绳机器人，
末端平台的广义力写为

\begin{equation}
  \bm W_p=
  \begin{bmatrix}
    F_x & F_y & M_z
  \end{bmatrix}^{\mathrm T},
\end{equation}

其中，$F_x$ 和 $F_y$ 分
别表示末端平台在 $x$
和 $y$ 方向上的等效广
义力，$M_z$ 表示绕平面
法向的等效广义力矩。

为了利用位姿误差实现
末端闭环补偿，本文采
用比例—微分—积分型
任务空间补偿器构造广
义力修正量

\begin{equation}
  \Delta \bm W_p
  =
  \bm K_p \bm e_q
  +
  \bm K_d \dot{\bm e}_q
  +
  \bm K_i \int \bm e_q\,\mathrm dt,
  \label{eq:ch4-deltaW}
\end{equation}

其中，
$\bm K_p$、$\bm K_d$、
$\bm K_i$ 分别为任务空
间比例、微分和积分增
益矩阵，通常取为对角
矩阵形式，即

\begin{equation}
  \bm K_p=
  \operatorname{diag}(k_{px},k_{py},k_{p\theta}),
\end{equation}

\begin{equation}
  \bm K_d=
  \operatorname{diag}(k_{dx},k_{dy},k_{d\theta}),
\end{equation}

\begin{equation}
  \bm K_i=
  \operatorname{diag}(k_{ix},k_{iy},k_{i\theta}).
\end{equation}

由此，位姿误差被转化
为对末端平台所需广义
力的修正需求。若同时
考虑轨迹规划或动力学
前馈项，则末端期望广
义力可进一步写为

\begin{equation}
  \bm W_p^\ast
  =
  \bm W_{ff}
  +
  \Delta \bm W_p,
  \label{eq:ch4-totalW}
\end{equation}

其中，$\bm W_{ff}$ 为基
于轨迹或动力学模型计
算得到的前馈项。在本
章中，重点放在闭环补
偿结构本身，因此前馈
项的具体构造将在后续
动力学与张力分配章节
中进一步讨论。

\begin{figure}[htbp]
  \centering
  \fbox{\begin{minipage}[c][4.8cm][c]{0.84\textwidth}
    \centering
    占位图：末端位姿闭环补偿示意图。\\
    建议给出末端期望位姿、正运动学反馈位姿、\\
    误差计算与广义力修正模块之间的关系。
  \end{minipage}}
  \caption{末端位姿闭环补偿示意图}
  \label{fig:ch4-task-space-loop}
\end{figure}

\subsection{广义力修正到各电机张力增量的映射}

根据第 3 章推导的结构
矩阵形式，四根绳张力
与末端平台广义力满足

\begin{equation}
  \bm W_p=\bm W \bm f,
  \label{eq:ch4-wf}
\end{equation}

其中，

\begin{equation}
  \bm f=
  \begin{bmatrix}
    f_1 & f_2 & f_3 & f_4
  \end{bmatrix}^{\mathrm T}
\end{equation}

为四根绳张力向量，
$\bm W\in\mathbb R^{3\times 4}$
为结构矩阵。由于系统
属于 4 绳驱动 3 自由度
形式，因此张力解具有
一维冗余。

设预紧张力向量为

\begin{equation}
  \bm f_0=
  \begin{bmatrix}
    f_{01} & f_{02} & f_{03} & f_{04}
  \end{bmatrix}^{\mathrm T},
\end{equation}

则在外环补偿作用下，
目标张力可写为

\begin{equation}
  \bm f^\ast
  =
  \bm f_0+\Delta \bm f.
  \label{eq:ch4-fstar}
\end{equation}

其中，张力增量 $\Delta \bm f$
由末端广义力修正量映
射得到。若暂不考虑零
空间调节，则可用广义
逆形式写为

\begin{equation}
  \Delta \bm f
  =
  \bm W^\dagger
  \Delta \bm W_p.
  \label{eq:ch4-deltaf-basic}
\end{equation}

进一步考虑冗余自由度
带来的张力整形能力，
则更一般的形式为

\begin{equation}
  \Delta \bm f
  =
  \bm W^\dagger
  \Delta \bm W_p
  +
  \left(
  \bm I-\bm W^\dagger\bm W
  \right)\bm z,
  \label{eq:ch4-deltaf-null}
\end{equation}

其中，$\bm z$ 为零空间调
节向量，用于后续实现
张力均衡、避免绳索松
弛以及限制峰值张力等
目标。由于本文下一章
将专门讨论冗余驱动下
的张力分配问题，因此
本章仅从控制实现角度
说明张力增量的基本映
射关系，而不展开最优
分配方法的细节。

当得到各绳目标张力后，
即可根据卷筒半径将其
转换为电机目标力矩

\begin{equation}
  \bm \tau^\ast
  =
  r_m \bm f^\ast
  =
  r_m
  \begin{bmatrix}
    f_1^\ast & f_2^\ast & f_3^\ast & f_4^\ast
  \end{bmatrix}^{\mathrm T}.
  \label{eq:ch4-tau-star}
\end{equation}

于是，末端位姿误差经
“广义力修正—张力增
量映射—电机目标力矩”
这一路径，最终转化为
四个单电机内环可执行
的控制参考量。

\begin{figure}[htbp]
  \centering
  \fbox{\begin{minipage}[c][5.0cm][c]{0.84\textwidth}
    \centering
    占位图：广义力修正到张力增量映射示意图。\\
    建议表示任务空间误差、广义力修正、\\
    张力映射、目标张力与目标力矩之间的层级关系。
  \end{minipage}}
  \caption{广义力修正到张力增量映射关系示意图}
  \label{fig:ch4-force-to-tension}
\end{figure}

\section{控制算法实现流程}

综合上述分析，本文控
制算法的整体实现过程
可总结为如下步骤。

第一步，根据任务要求
生成末端参考轨迹，得
到期望位姿、期望速度
以及必要的前馈信息。

第二步，通过逆运动学
和正运动学模块建立任
务空间与执行空间之间
的信息通道，实时获得
末端平台位姿估计值。

第三步，根据期望位姿
与反馈位姿计算末端误
差，并利用任务空间补
偿器构造广义力修正量。

第四步，根据广义力与
绳张力之间的映射关系，
将末端广义力修正需求
转换为四根绳的目标张
力及目标力矩。

第五步，对四个执行单
元分别运行 PID、摩擦
补偿和扰动观测器构成
的单电机内环控制器，
实现各电机目标力矩跟
踪。

第六步，在实时控制循
环中不断重复上述过程，
直至完成整个运动任务。

\begin{figure}[htbp]
  \centering
  \fbox{\begin{minipage}[c][5.0cm][c]{0.84\textwidth}
    \centering
    占位图：控制算法实现流程图。\\
    建议按“参考轨迹生成—位姿反馈—误差计算—\\
    广义力修正—张力映射—单电机内环控制”顺序绘制。
  \end{minipage}}
  \caption{控制算法实现流程图}
  \label{fig:ch4-algorithm-flow}
\end{figure}

为保证算法在样机系统
中的可实现性，控制器
设计中遵循以下原则：
一是内环结构尽量简洁，
便于参数整定和实时实
现；二是外环误差补偿
采用任务空间形式，便
于与末端轨迹目标直接
对应；三是张力映射保
持与后续冗余张力分配
分析一致，以便形成从
控制器到实验平台的闭
环实现链条。基于上述
设计思想，本文后续实
验将分别从单电机力矩
跟踪、四电机协同张力
控制以及末端加速度跟
踪三个层面，对控制效
果进行验证。

\section{本章小结}

本章围绕样机系统的控
制实现需求，建立了单
电机闭环力矩控制与末
端位姿闭环补偿相结合
的控制算法框架。首先，
从控制系统总体架构出
发，明确了“单电机内环
力矩控制+末端外环位姿
补偿”的双层控制思路；
随后，针对单电机执行
单元建立了基于 PID、
摩擦补偿和扰动观测器
的闭环力矩控制方法，
提高了执行单元的跟踪
精度与抗扰能力；进一
步地，利用正运动学反
馈构建末端位姿闭环补
偿方法，将末端轨迹误
差转化为平台广义力修
正量，并建立了广义力
修正到各绳张力增量及
电机目标力矩之间的映
射关系。

本章所建立的控制算法
为后续冗余驱动下的张
力分配与四电机协同控
制研究提供了实现基础。
在此基础上，下一章将
进一步结合动力学模型，
围绕张力分配策略、可
行张力区间及不同分配
方法的效果展开分析，
从而形成更完整的协同
控制方案。